\section{结果分析}
\subsection{输入视图数与深度预测结果}

\subsubsection{深度预测对比}

保持其他设置不变，分别使用 3、7、11 张输入图像进行深度预测。图中选取 0012、0025 和 0043 三个参考视角，比较建筑主体、轮廓和背景区域的变化；表中统计使用全部 49 个参考视角。

\ReportFigure[307bp]{depth_comparison.png}{不同输入视图数的深度预测。由左至右为参考图、3、7、11 视图深度和 7 视图置信度。深度使用同一色阶，置信度图中白色表示较高、黑色表示较低，显示范围截取为 0–1。}


\subsubsection{视图数变化的影响}
\input{view-comparison-table.tex}
\ReportNote{最终保留比例与融合点数均采用置信度阈值 0.8；平均置信度和比例按 49 张等尺寸深度图统计。}

三组深度图都能区分建筑的前后关系，主体形状比较接近。增加输入视图后，较明显的变化出现在建筑轮廓和背景区域，同一位置在不同组中会出现不同的深度颜色。部分区域有所变化，但从整幅图看，没有出现输入视图越多、各处细节就越清楚的趋势。

表 2 中，3、7 视图的平均置信度和点云保留比例也比较接近，11 视图的保留比例则略低。采用相同后处理后，三组融合点数分别约为 99.8、99.4、95.2 万。结合图像和统计结果，本次增加输入视图主要改变了局部预测，没有带来整体结果的持续改善。因此，比较输入视图数时，还需要看具体区域的变化，不能只凭图像数量判断效果。

